% part: intuitionistic-logic
% chapter: propositions-as-types
% section: types
%READERNOTE{SOL6-A319: प्रत्येक पदाला प्रकार असतो हा असत्य दावा काढला: योग्य प्रकार मिळाल्यास तो एकमेव असतो. सर्व संदर्भ-खुणांची सुसंगती आणि बद्ध चलांची सुरक्षित नावे स्पष्ट केली.}

\documentclass[../../../include/open-logic-section]{subfiles}

\begin{document}

\olfileid{int}{pty}{typ}

\olsection{प्रकार}

$(MN)$ सारखे सिद्धता-पद स्वतःहून कोणत्या सिद्धतेचे
सिद्धता-पद आहे हे एकमेवपणे ठरवत नाही. कारण मुक्त चले,
म्हणजे स्वयंसिद्धकांची सिद्धता-पदे, संबंधित स्वयंसिद्धकांतील
!!{formula}s सांगत नाहीत. पण ती सुसंगत संदर्भ~$\Gamma$ मध्ये
दिली आणि पद योग्य प्रकाराचे असेल, तर सिद्धता \emph{एकमेवपणे} ठरते.

\begin{defn}
सिद्धता-पद~$N$ साठी \emph{संदर्भ} म्हणजे प्रकार-खुणांचा
सुसंगत संच~$\Gamma$: कोणत्याही चलाला दोन वेगळे प्रकार
देता येत नाहीत, आणि~$N$ च्या प्रत्येक मुक्त चलासाठी
$\Gamma$ मध्ये $x:!A$ अशी नेमकी एक जोडी असते.
\end{defn}

या व्याख्येनुसार~$\Gamma$ मध्ये $x:!A$ ही जोडी
\emph{फक्त}~$N$ च्या मुक्त चलांसाठीच असावी अशी अट नाही.
उदाहरणार्थ, $\Gamma = \{x:!A, y:!B, z:!C\}$ हा
$\pair{x}{y}$ साठी संदर्भ आहे.

\begin{defn}\ollabel{defn:type}
  ~$\Gamma$ हा~$N$ साठी संदर्भ आहे असे माना.
  ~$\Gamma$ च्या सापेक्ष~$N$ चा \emph{प्रकार}
  पुढीलप्रमाणे विगमनाने परिभाषित करतो:
  \begin{enumerate}
  \item $x:!A \in \Gamma$ असेल तेव्हाच~$x$ चा प्रकार~$!A$ असतो.
  \item $x:!A, \Gamma$ च्या सापेक्ष~$N$ चा प्रकार~$!B$
    असेल, तर $\lambd[\typeof{x}{!A}][N]$ चा प्रकार
    $!A \lif !B$ असतो.
  \item $N$ चा प्रकार $!A \lif !B$ आणि~$M$ चा प्रकार~$!A$
    असेल, तर $(NM)$ चा प्रकार~$!B$ असतो.
  \item $N$ चा प्रकार~$!A$ आणि~$M$ चा प्रकार~$!B$
    असेल, तर $\pair{N}{M}$ चा प्रकार $!A \land !B$ असतो.
  \item $N$ चा प्रकार~$!A_1 \land !A_2$ असेल, तर
    $\proj{i}{N}$ चा प्रकार~$!A_i$ असतो.
  \item $N$ चा प्रकार~$!A$ असेल, तर
    $\inj[!B]{i}{N}$ चा प्रकार $i=1$ असताना
    $!A \lor !B$ आणि $i=2$ असताना $!B \lor !A$ असतो.
  \item $M$ चा प्रकार $!A \lor !B$, तसेच
    $x_1:!A, \Gamma$ च्या सापेक्ष~$N_1$ चा प्रकार~$!C$
    आणि $x_2:!B, \Gamma$ च्या सापेक्ष~$N_2$ चा
    प्रकार~$!C$ असेल, तर
    $\dcase{M}{x_1}{N_1}{x_2}{N_2}$ चा प्रकार~$!C$ असतो.
  \item $N$ चा प्रकार~$\lfalse$ असेल, तर
    $\abort{!A}{N}$ चा प्रकार~$!A$ असतो.
  \end{enumerate}
\end{defn}
या नियमांत न बसणाऱ्या पदाला प्रकार नाही. अमूर्तीकरणाच्या
आणि प्रकरण-विभाजनाच्या बद्ध चलांना संदर्भातील चलांपासून वेगळी
नवी नावे देऊन वरील संदर्भ-वाढ करायची.

\begin{prop}
सिद्धता-पद~$N$ साठीच्या सुसंगत संदर्भ~$\Gamma$ च्या
सापेक्ष~$N$ ला जास्तीत जास्त एक प्रकार असतो. प्रकार
अस्तित्वात असेल, तर तो एकमेव असतो; प्रत्येक कच्च्या
सिद्धता-पदाला प्रकार मिळतोच असे नाही.
\end{prop}

\begin{proof}
  ~$N$ वर विगमनाने.
\end{proof}

संदर्भ~$\Gamma$ च्या सापेक्ष~$N$ चा प्रकार~$!A$
असेल, तर $\Gamma \Sequent N : !A$ असे लिहितो.
\olref{defn:type} मधील उपवाक्ये परिचित वाटली पाहिजेत:
त्यात एक लहानसा फरक वगळता \Log{tN2i} चेच नियम
आहेत—संदर्भ~$\Gamma$ सर्वत्र तोच राहतो. ही
व्याख्या \Log{tN3i} या कलनाच्या रूपात मांडता येते;
त्याचे नियम \olref{tab:tN3ip} मध्ये दिले आहेत.

\subfile{rules-tN3}

सिद्धता-पदे ही $\lambd$-कलनाच्या एका आवृत्तीतील
पदे आहेत: \emph{गुणाकार, बेरीज आणि रिक्त प्रकार
असलेले प्रकारयुक्त $\lambd$-कलन}. पारंपरिक
$\lambd$-कलनात फक्त चलांपासून बनलेली पदे,
\emph{$\lambd$-अमूर्तीकरणे}~$\lambd[x][N]$, आणि
उपयोजन $(MN)$ असते; त्यात प्रकारयुक्त
$\lambd$-अमूर्तीकरण~$\lambd[\typeof{x}{!A}][N]$
यातील~$!A$ हे प्रकार-चिन्ह नसते. ~$\lambd$-कलन
हे ट्यूरिंग-संपूर्ण संगणन-प्रतिमान आहे (म्हणजे प्रत्येक
आंशिक संगणनीय फलन त्यातील पदाने दर्शवता येते).
सिद्धता-पदे आणि प्रकार-नियम देतात ती~$\lambd$-कलनाची
आवृत्ती निर्बंधित संगणन-प्रतिमान आहे; पण तिच्या मूलभूत
कल्पना त्याच आहेत. हे निर्बंध प्रकार-निर्देशनातून येतात:
फक्त \emph{योग्य प्रकाराची} पदे, म्हणजे
संदर्भ~$\Gamma$ च्या सापेक्ष प्रकार असलेली
सिद्धता-पदे, ग्राह्य असतात. उदाहरणार्थ,
$\lambd[x][(xx)]$ हे पारंपरिक प्रकारविरहित
$\lambd$-कलनात ग्राह्य पद आहे; पण
$\lambd[\typeof{x}{!A}][(xx)]$ हे कोणत्याही~$!A$
साठी योग्य प्रकाराचे नाही. कारण $(xx)$ ला प्रकार
असण्यासाठी पहिल्या~$x$ चा प्रकार $!A \lif !B$
आणि दुसऱ्या~$x$ चा प्रकार~$!A$ असावा लागेल.
पण~$!A$ हे $!A \lif !B$ चे योग्य उप-!!{formula}
असल्यामुळे हे अशक्य आहे.

एखाद्या प्रकाराचे पद ज्या वस्तूला नाव देते, तिचा
प्रकार म्हणून प्रकारांकडे सहजपणे पाहता येते. म्हणून
$!A \land !B$ या प्रकाराचे पद पहिला घटक~$!A$
आणि दुसरा~$!B$ या प्रकाराचा असलेली जोडी निवडते;
म्हणजे गुणाकार~$!A \times !B$ मधील !!a{element}.
$!A \lif !B$ या प्रकाराचे पद प्रकार~$!A$ च्या
वस्तूंपासून प्रकार~$!B$ च्या वस्तूंकडे जाणारे फलन आहे.
$!A \lor !B$ या प्रकारचे पद हे प्रकार~$!A$ किंवा
प्रकार~$!B$ ची वस्तू असते; तसेच ती दोहोंपैकी कोणत्या
प्रकाराची आहे ही माहिती त्यात असते. हे~$!A$ आणि~$!B$
या दोन संचांच्या असंयुक्त संघासारखे आहे:
$\{\tuple{i,x} : i = 1 \text{ किंवा } 2, x \in !A
\text{ जर } i = 1, x \in !B \text{ जर } i=2\}$.
$\lfalse$ हा प्रकार रिक्त संच आहे; म्हणजे त्या प्रकाराचे
पद कोणत्याही वस्तूला नाव देऊ शकत नाही.
नैसर्गिक निगमन (\Log{tN3i} सह) सुसंगत असल्यामुळे
खुली गृहीतके नसताना~$\lfalse$ ची !!{proof} नसते.
म्हणून~$\lfalse$ प्रकाराचे बंद सिद्धता-पद
(मुक्त चले नसलेले पद) नसते.

प्रकारयुक्त $\lambd$-कलनाचे संकारक योग्य प्रकाराच्या
वस्तूंना सहजपणे नाव देणारी पदे तयार करतात, अर्थात
त्यांना लावलेली पदेही ठरावीक प्रकारच्या वस्तूंना नाव
देत असतील तर. उदाहरणार्थ, ``$N_1 : !A$ आणि
$N_2 : !B$ असतील, तर
$\pair{N_1}{N_2} : !A \land !B$'' याचा अर्थ:
$N_1$ ने प्रकार~$!A$ च्या वस्तूला आणि~$N_2$ ने
प्रकार~$!B$ च्या वस्तूला नाव दिले, तर
$\pair{N_1}{N_2}$ हे प्रकार $!A \land !B$ च्या
वस्तूला नाव देते.

\end{document}
